\chapter{Sign-off}
\label{ch:signoff}

The \term{sign-off} is the last step of a run. It decides whether hdl-vet writes the
output tree. You see it when it fails, so this chapter is about what each failure
says and what to do about it.

\section{What the sign-off does}

When all scripts have run, hdl-vet does these steps.

\begin{enumerate}
\item It applies the queued edits to the text.
\item It adds the header line \code{// Modified by hdl-vet; scripts: ...} to every file
  that differs from its original. The header is part of the text that is checked.
  What is checked is byte for byte what would be written.
\item If a file differs, it loads the edited design and compares it with the design
  that you started from, the \term{baseline}.
\item It evaluates the checks that the operations registered.
\end{enumerate}

If every step passes, the output tree is written. If a file did not change, there is
nothing to compare, and the checks run on the baseline. Everything is compared
against the baseline, not against zero: a warning that the generator's design
already has is not your problem, and does not fail the sign-off.

A failed sign-off prints \code{hdl-vet sign-off: FAILED} and one line for each
problem, to the error output. The exit status is 1. No output tree is written. This
holds with \opt{--dry-run} too.

The exit status is 1 when the sign-off fails. Chapter~\ref{ch:cli} lists all exit statuses.

\section{Every line of the report}

Each problem is one line. The lines come in this order. The transcript after the list
shows one failure of each kind.

\begin{description}
\item[the edited design cannot be loaded: REASON] The compiler cannot read the
  edited text, or cannot find the top module. The reason is its message. Look at
  your raw edits, or at an operation that changed a name that the top module needs.
  An operation that fails in this way alone is usually refused before.

\item[new file read by the edited design, not in the baseline: PATH] The edited
  design reads a file that the original design did not read. A raw edit added an
  \code{`include}. Use \apiname{Vet.add_file} to bring in a file on purpose.

\item[new diagnostic: SEVERITY CODE at FILE:LINE: MESSAGE] The compiler reports an
  error or a warning that the baseline did not have. \code{CODE} is the compiler's
  name for it, for example \code{PortWidthExpand}. \code{FILE} and \code{LINE} are
  in the edited text. The compiler reports its errors and its default warnings,
  and hdl-vet switches on five more that an edit of ports can cause:
  \code{ImplicitNet}, \code{PortWidthTruncate}, \code{PortWidthExpand},
  \code{WidthTruncate} and \code{EmptyInputPortConn}. Notes of the compiler are
  not counted. A diagnostic that is in the
  baseline and has only moved because of an edit is not new. Warnings that a file
  already had when a script added it are not new, but a later warning located in that file
  counts. Read the message and the line, and fix the
  script.

\item[left dangling: MESSAGE in module MODULE (FILE:LINE)] Something is read but
  nothing drives it. See the next section.

\item[unexpected change: KIND PATH NAME: BEFORE -> AFTER] The structure of the
  design changed in a way that no operation and no \code{expect} or \code{permit}
  explains. \code{KIND} is the kind of change (see \apiname{Vet.expect}), \code{PATH}
  the instance, \code{NAME} the port or parameter, if any. The values are the text
  before and after. A change that you made by hand needs an \code{expect}. A change
  that you did not expect from an operation shows that the operation touched more than
  you thought.

\item[expected change not seen: KIND path=PATH name=NAME] Something was expected
  (\code{expect}, or an operation that must change the structure) and the design
  does not show it. Another edit may have undone it, or the \code{expect} is wrong.

\item[failed check: MESSAGE] A post-condition failed. See the section after next.
\end{description}

These scripts produce the failures in the transcript. \file{cannot_load.py} renames the
top module with a raw edit. \file{new_file.py} inserts an \code{`include}.
\file{wide_port.py} connects a 1-bit expression to a 2-bit port. \file{dangling.py}
opens an output that something reads. \file{unexpected.py}, \file{unseen.py} and
\file{add_cell.py} are in chapter~\ref{ch:operations}.

\lstinputlisting[style=script]{examples/scripts/cannot_load.py}
\lstinputlisting[style=script]{examples/scripts/new_file.py}
\lstinputlisting[style=script]{examples/scripts/wide_port.py}
\lstinputlisting[style=script]{examples/scripts/dangling.py}

\lstinputlisting[style=console]{examples/console/signoff.txt}

The last-but-one command runs two scripts that fail in different ways. The report has
one line for each problem.

\section{Dangling: an output may float, an input may not}

The sign-off compares the \term{dangling set} of the final design with that of the
baseline. Only new entries fail. The set holds:

\begin{itemize}
\item a net or variable that is read but never driven;
\item an output port that nothing inside its module drives, \emph{but only when
  something reads it}: some instance of the module connects it to a non-empty
  expression, or the module itself reads its signal.
\end{itemize}

An output that nobody reads may float. It is allowed. This is the case when
every instance leaves it open (\code{.p()}), and for an output of the top module.
An input may not float. A new input port that no instance connects is a new
diagnostic. An input connected to nothing with \code{.p()} is a new diagnostic too.
Entries inside an added file never count.

The transcript shows a run in which \file{dangling.py} disconnects the output
\code{irq_o} of \code{u_periph}. The signal \code{irq} is read by the top module
and is now driven by nobody: \emph{left dangling}. In chapter~\ref{ch:operations},
\file{ports.py} disconnects \code{count1_o}, which goes to an output of the top
module: that is allowed.

The sign-off does not see everything. Chapter~\ref{ch:limits} lists what it misses,
for example a signal that loses the driver of only some of its bits.

\section{Post-conditions}

Each applied operation registers a \term{check} that is evaluated on the final
design. The check makes sure that a later edit has not undone the operation. A failed
check starts with the name of the operation and the target. Examples:

\begin{itemize}
\item \code{failed check: add_wire soc_top tmp: declaration not found in the final
  design}
\item \code{failed check: connect soc_top.u_rst_sync test_i: expected 'x', found 'y'}
\item \code{failed check: add_file NAME: none of its modules is instantiated}
\item \code{failed check: propagate_param WIDTH PATH: the value in use changed
  from 8 to 4}
\end{itemize}

If you see one, look at the later operations of your script (or of a later script):
one of them has changed what this operation wrote. Two operations that cancel each
other, such as \apiname{Module.add_wire} and \apiname{Module.remove_wire} on the same name,
fail in this way.

Skipped operations register no check.

\section{Expectations}

The sign-off compares the structure of the baseline and the final design, instance by
instance. A \term{change} has a kind, the path of the instance, a name and the text
before and after.

\begin{description}
\item[instance-added, instance-removed] An instance appears or goes. An added or
  removed instance shows only this change, and the same change for each instance
  below it.
\item[module-changed] The instance uses another module.
\item[param-changed] A parameter has another value at the instance.
\item[port-added, port-removed, port-changed] The module of the instance gains, loses or
  changes a port. A port is described by its direction and width, for example
  \code{input 1}.
\item[port-type-changed] The direction and width stay, but the declared type differs.
\item[connection-changed] The connection of a port has other text.
\end{description}

Every change must be covered by an expectation or a permission, and every
expectation must match a change. Operations register their own. For example,
\apiname{Module.add_port} requires \code{port-added} at every path of the module, and
permits \code{connection-changed} for the port. You add yours with \apiname{Vet.expect}
and \apiname{Vet.permit}.

\section{Name the top module}

hdl-vet needs to know the top module of your design. If you give none with
\opt{--top}, it takes the modules that nobody instantiates. There is a trap. If your
script replaces or removes the \emph{last} instance of a module, that module is no
longer instantiated, so it becomes a top-level module of the edited design. The
sign-off then reports \code{instance-added}: a new top appeared. This is not a
problem of your script. Name the top module, and the edited design has the same top
as the original. The transcript shows both. Put \code{top} in the project file so
that nobody forgets.

\lstinputlisting[style=console]{examples/console/top.txt}
